% =============================================================================
%  وثيقة التلميذ — الأعمال التطبيقية رقم 01 (Fiche Élève TP N° 01)
%  المجال 1: بيئة التعامل مع الحاسوب — الوحدة 1: تقنية المعلومات (IT)
%  عنوان الحصة: استكشاف مواصفات الحاسوب وتجريب دورة معالجة البيانات
%  السنة الدراسية: 2026 / 2027
% =============================================================================
%  البناء (من هذا المجلد):
%    xelatex -interaction=nonstopmode -halt-on-error -output-directory=build student-sheet.tex
% =============================================================================

\documentclass[11pt, a4paper]{article}

% ── المشترك ─────────────────────────────────────────────────────────────────
\input{"../../shared/common.tex"}

% ══════════════════════════════════════════════════════════════════════════════
\begin{document}

% ── الصفحة 1: الترويسة + العنوان + بيانات التلميذ ───────────────────────────
\officialheader

\sheettitle{01}{استكشاف مواصفات الحاسوب وتجريب دورة معالجة البيانات}
           {مخبر المعلوماتية}{01 ساعة}

\studentid

% ── الصفحة 2: المهمة الأولى ────────────────────────────────────────────────
\sectionbreak{المهمة الأولى: استكشاف مواصفات العتاد ونظام التشغيل}

\begin{activitybox}[التعليمات]
  قم بالضغط بالزر الأيمن للفأرة على أيقونة «هذا الكمبيوتر» (Ce PC / This PC)، ثم اختر «خصائص» (Propriétés)، ودوّن مواصفات جهازك في الجدول التالي:
  \vspace{4pt}
  \specheader
  \specrow{نوع وسرعة المعالج (Processor / CPU)}{1}
  \specrow{سعة الذاكرة الحية (RAM)}{1}
  \specrow{اسم ونوع نظام التشغيل (OS)}{1}
  \specrow{سعة القرص الصلب المتاح (D:)}{1}
  \specend
\end{activitybox}

% ── الصفحة 3: المهمة الثانية ────────────────────────────────────────────────
\sectionbreak{المهمة الثانية: تجسيد دورة معالجة البيانات}

\begin{activitybox}[التعليمات]
  \textbf{1.} افتح تطبيق محرّر النصوص أو المجدول المتاح على جهازك.\newline
  \textbf{2.} \textbf{إدخال البيانات الخام (Data):} قم بكتابة قائمة التلاميذ التالية كما هي دون تعديل:
  \begin{center}
  \begin{tabular}{@{}c@{}}
    سارة — 2009 — 09.50 \\
    أحمد — 2008 — 14.00 \\
    مريم — 2008 — 11.50 \\
    كريم — 2009 — 08.00 \\
  \end{tabular}
  \end{center}
  \textbf{3.} \textbf{إجراء المعالجة (Processing):}
  \begin{itemize}[label=\textbullet, leftmargin=1.1em, itemsep=1pt, topsep=2pt, parsep=0pt]
    \item قم بإعادة ترتيب الأسماء أبجديًا من الألف إلى الياء.
    \item أضف خانة «الملاحظة» أمام كل تلميذ حسب القاعدة: إذا كان المعدل $\ge 10.00$ $\rightarrow$ \textbf{ناجح}، وإذا كان $< 10.00$ $\rightarrow$ \textbf{مؤجل}.
  \end{itemize}
  \textbf{4.} \textbf{استخراج المعلومات (Information):} سجّل النتائج النهائية بعد المعالجة:
  \begin{itemize}[label=\textbullet, leftmargin=1.1em, itemsep=3pt, topsep=2pt, parsep=0pt]
    \item عدد الناجحين: \blank[2.2cm]
    \item عدد المؤجلين: \blank[2.2cm]
    \item التلميذ المتحصّل على أعلى معدل: \blank[5cm]
  \end{itemize}
\end{activitybox}

% ── الصفحة 4: المهمة الثالثة ────────────────────────────────────────────────
\sectionbreak{المهمة الثالثة: تنظيم وتخزين الملف}

\begin{activitybox}[التعليمات]
  \begin{enumerate}[label=\arabic*., leftmargin=1.4em, itemsep=2pt, topsep=1pt, parsep=0pt]
    \item افتح القرص الصلب (D:).
    \item أنشئ مجلدًا جديدًا واجعله باسمك وفق الشكل التالي: \textbf{اسمك\_لقبك\_الفوج}.
    \item احفظ ملف عملك داخل مجلدك باسم: \texttt{TP01\_Informatique}.
  \end{enumerate}
  \vspace{4pt}
  اكتب المسار الكامل لحفظ ملفك هنا: \\[2pt]
  \noindent D:\textbackslash\ \blank[12cm]
\end{activitybox}

% ── الصفحة الأخيرة: جدول التقييم (خاص بالأستاذ) ─────────────────────────────
\sectionbreak{جدول تقييم عمل التلميذ (خاص بالأستاذ)}

{\gridsetup
\centering
\begin{tabularx}{0.9\textwidth}{@{}|X|>{\centering\arraybackslash}p{2.8cm}|@{}}
  \hline
  \rowcolor{primary}
  \thead{المعيار} & \thead{الملاحظة والدرجة} \\
  \hline
  \graderow{1. التشغيل والاستعمال الآمن للجهاز والالتزام بالنظام.}{/ 04}
  \graderow{2. استخراج مواصفات الجهاز بدقة (المهمة 1).}{/ 05}
  \graderow{3. إدخال البيانات ومعالجتها واستخراج المعلومة الصحيحة (المهمة 2).}{/ 06}
  \graderow{4. إنشاء المجلد وحفظ الملف في المسار الصحيح (المهمة 3).}{/ 05}
  \graderow{\textbf{المجموع النهائي}}{/ 20}
\end{tabularx}\par
}\vspace{4pt}

\end{document}
